\chapter*{Preface}

\textbf{About These Notes:} These notes aim to introduce the concept of cofinality, offering the motivations and intuitions behind it. The initial chapters consist primarily in results and examples that explore the idea of cofinality, showing how it naturally arises and examining its key properties.

From there on, almost everything presented can be found in Jech's book - Introduction to Set Theory, all I will do here is a comprehensive and filtered summary, along with some of my own intuitions that may be useful to someone, and some additional comments on things I've found elsewhere. The full project can be found in my \href{https://github.com/Xennonio/Cofinality/tree/main}{Github Repository}

\textbf{Prerequisites:} It is important to note that these notes assume the reader has some familiarity with the basics of cardinal and ordinal arithmetic to fully understand the proofs of the theorems. However, even without this background, readers can still grasp the underlying intuition by skipping the proofs and following the comments and ramblings throughout the text, which, hopefully, will at least guides you through a natural exploration of the defined concepts.

\textbf{Further Directions:} As you will see in the notes, the concept of cofinality leads us directly to Large Cardinal Axioms by addressing a natural question in the study of regular and singular cardinals: "does there exist an uncountable regular limit cardinal?", from here, you can read my notes on Large Cardinals (which I'm still working on). Furthermore, cofinality also naturally leads us to cardinal arithmetic, which is full of consistency results that, if you are interested, you can read more about in my notes on \href{https://github.com/Xennonio/Consistency-Results}{Consistency Results} (which I'm also still working on).